\section{Reinforcement Learning from Human Feedback}
\label{sec:rlhf}

The structural estimation and IRL literatures are primarily concerned with recovering an agent's underlying reward function from observations of their enacted, often optimal, behavior. A distinct but related paradigm, Reinforcement Learning from Human Feedback (RLHF), has recently emerged to address settings where expert demonstrations are either unavailable, prohibitively expensive to collect, or may not represent the true desired behavior. Instead of inverting an observed policy $\pi$ to find the reward function that rationalizes it, RLHF learns a reward function directly by eliciting preferences from a human evaluator over pairs of machine-generated behaviors. While this approach has become the predominant method for aligning large language models \citep{ouyang2022training}, its methodology is general and was first demonstrated as a viable technique for training agents in standard reinforcement learning domains, such as simulated robotics and video games, without access to an explicit, pre-defined reward function \citep{christiano:2017}.

The canonical RLHF framework is best understood as a three-step process that combines supervised learning, preference modeling, and policy optimization. The objective is to fine-tune a base policy, typically a large pretrained model, to align its behavior with human judgments. The first step, supervised fine-tuning (SFT), initializes the policy in a sensible region of the parameter space. A small, high-quality dataset of human-written demonstrations of the desired task is collected, and a pretrained policy is fine-tuned via standard maximum likelihood to produce an initial policy, which we denote $\pi^{SFT}$.

The second and central step is the learning of a reward model, $r_\theta(s, y)$, parameterized by $\theta$. Note that unlike the per-action rewards $r(s,a)$ in the preceding MDP framework, the reward here is assigned to a complete trajectory or output $y$ generated from state $s$. This step directly applies the principles of discrete choice modeling discussed in Section \ref{sec:ddc}. A dataset $\mathcal{D}$ of human preferences is collected, consisting of tuples $(s, y_w, y_l)$, where $s$ is a context or state, and $y_w$ and $y_l$ are two trajectories or outputs generated by the policy, with $y_w$ being the "winner" preferred by the human evaluator and $y_l$ the "loser." The reward model $r_\theta$ is then trained to predict these choices. Assuming the human preferences can be described by a random utility model, where the latent utility of an output $y$ in context $s$ is given by $r_\theta(s, y)$, the probability that $y_w$ is preferred over $y_l$ can be modeled using a logit or softmax function. The reward model parameters $\theta$ are estimated by minimizing the negative log-likelihood of the human choices, which corresponds to a binary cross-entropy loss:
\begin{equation}
    \mathcal{L}(\theta) = -\mathbb{E}_{(s, y_w, y_l) \sim \mathcal{D}} \left[ \log \sigma \left( r_\theta(s, y_w) - r_\theta(s, y_l) \right) \right],
\label{eq:rlhf_loss}
\end{equation}
where $\sigma(\cdot)$ is the logistic sigmoid function. The learned $r_\theta$ thus serves as a scalar proxy for human preferences.

In the third and final step, the policy is optimized using reinforcement learning, with the learned reward model $r_\theta$ providing the reward signal. The policy, initialized from $\pi^{SFT}$, is updated to maximize the expected reward given by $r_\theta$. A critical component of this optimization, first introduced in this context by \citet{ziegler2019fine}, is the inclusion of a regularization term that penalizes the divergence of the learned policy, $\pi_\phi$, from the initial supervised policy, $\pi^{SFT}$. The policy is then optimized to maximize the expected reward from the human-preference model, regularized by a KL-divergence penalty against the initial policy $\pi^{SFT}$. The objective function for the policy optimization is:
\begin{equation}
    J(\phi) = \mathbb{E}_{s \sim \mathcal{D}, y \sim \pi_\phi(\cdot|s)} [r_\theta(s, y)] - \beta_{KL} \mathbb{E}_{s \sim \mathcal{D}} [D_{KL}(\pi_\phi(\cdot|s) \,||\, \pi^{SFT}(\cdot|s))],
\label{eq:rlhf_objective}
\end{equation}
where $D_{KL}$ is the Kullback-Leibler divergence and $\beta_{KL}$ is a hyperparameter controlling the strength of the penalty. The KL-divergence term in this objective serves two essential functions: it acts as an entropy bonus to encourage exploration and, more importantly, it constrains the policy to remain close to the distribution of coherent outputs learned during the SFT phase. This prevents the policy from learning to generate adversarial or out-of-distribution outputs that "hack" the fixed reward model by exploiting its imperfections to achieve a high score with degenerate behavior.

While RLHF has proven exceptionally effective at aligning model behavior with human preferences \citep{stiennon2020learning}, several open issues remain. The framework does not aim to solve the identification problem in the strict econometric sense; the learned reward model $r_\theta$ is a proxy for preference, not necessarily the uniquely identified "true" utility function required for welfare analysis. Furthermore, the notion of "human feedback" is a significant simplification, as the preferences being optimized are those of a small, non-representative group of paid labelers, raising important questions about whose values are being embedded in these systems. The practical implementation also presents a substantial operational challenge, as collecting a large, high-quality dataset of preferences requires a significant investment in both cost and meticulous quality control. Finally, the fine-tuning process can lead to a degradation of performance on standard academic benchmarks, a phenomenon sometimes referred to as an "alignment tax." Mitigating this requires further methodological refinements, such as mixing pretraining gradient updates into the policy optimization process, as demonstrated by the PPO-ptx algorithm of \citet{ouyang2022training}. Despite these challenges, RLHF represents a powerful and scalable approach for directing the behavior of complex generative agents in domains where rewards are difficult to specify but straightforward to evaluate through comparative judgment.

\subsection{Recent Developments: Direct Policy Optimization and Self-Rewarding}

Building upon the canonical RLHF framework, recent research has sought to simplify the multi-stage training pipeline and address the potential instabilities of reinforcement learning. A significant development in this direction is Direct Preference Optimization (DPO), a method introduced by \citet{rafailov2023direct}. DPO reframes the entire alignment problem as a direct supervised learning task, eliminating the need for an explicitly trained reward model and the complexities of RL-based policy optimization.

The key insight of DPO is the derivation of a closed-form relationship between the optimal policy of the KL-regularized reward maximization problem and the underlying reward function itself. For notational simplicity in the following derivation, we let $r(s,y)$ denote the ground-truth reward function that the trained model $r_\theta(s,y)$ aims to approximate. Recall that the theoretical optimal policy $\pi_r$ for the objective in equation (\ref{eq:rlhf_objective}) can be expressed in terms of a reference policy $\pi^{SFT}$ and the reward function $r(s,y)$ as:
\begin{equation}
    \pi_r(y|s) = \frac{1}{Z(s)} \pi^{SFT}(y|s) \exp\left(\frac{1}{\beta_{KL}} r(s,y)\right),
\label{eq:dpo_optimal_policy}
\end{equation}
where $Z(s)$ is a partition function that depends only on the state $s$. By rearranging this expression, the reward function can be reparameterized in terms of the optimal policy $\pi_r$ and the reference policy $\pi^{SFT}$:
\begin{equation}
    r(s,y) = \beta_{KL} \log\left(\frac{\pi_r(y|s)}{\pi^{SFT}(y|s)}\right) + \beta_{KL} \log(Z(s)).
\label{eq:dpo_reparameterize}
\end{equation}
Substituting this reparameterized reward into the preference model loss from equation (\ref{eq:rlhf_loss}) causes the state-dependent partition function term $\beta_{KL} \log(Z(s))$ to cancel out. This yields a loss function that depends only on the policy $\pi$ being optimized and the fixed reference policy $\pi^{SFT}$. The DPO loss is thus a simple binary cross-entropy objective over policy likelihoods:
\begin{equation}
    \mathcal{L}_{DPO}(\phi; \pi^{SFT}) = -\mathbb{E}_{(s, y_w, y_l) \sim \mathcal{D}} \left[ \log \sigma \left( \beta_{KL} \log \frac{\pi_\phi(y_w|s)}{\pi^{SFT}(y_w|s)} - \beta_{KL} \log \frac{\pi_\phi(y_l|s)}{\pi^{SFT}(y_l|s)} \right) \right].
\label{eq:dpo_loss}
\end{equation}
This objective can be optimized directly on the policy parameters $\phi$ using standard supervised learning techniques with a static, offline preference dataset. The gradient of this loss intuitively increases the likelihood of preferred responses $y_w$ while decreasing the likelihood of dispreferred responses $y_l$, relative to the reference policy. By recasting the constrained optimization problem as a simple likelihood maximization, DPO offers a more stable, computationally efficient, and easier-to-implement alternative to the multi-stage RLHF pipeline.

A subsequent innovation, Self-Rewarding Language Models, leverages the DPO framework to create an iterative self-improvement loop that reduces the reliance on static, human-annotated data \citep{yuan2024self}. In this paradigm, a single language model acts as both the instruction-following agent and its own reward model. The process begins with a base model that is fine-tuned on a seed dataset to learn both instruction-following and evaluation capabilities (LLM-as-a-Judge). In each subsequent iteration, the current model generates a new preference dataset for itself by producing multiple responses to a set of prompts and then evaluating its own outputs to assign scores. These scores are used to construct preference pairs $(y_w, y_l)$, which form an AI-generated feedback dataset. The model is then fine-tuned on this new dataset using the DPO loss from equation (\ref{eq:dpo_loss}). This iterative process, where the model generates its own training data and provides its own reward signal, has been shown to create a "virtuous circle" of improvement, where enhancements in instruction-following ability lead to better reward modeling, which in turn fuels the next round of policy improvement. This approach demonstrates a path toward scaling alignment beyond the limits of human data collection by enabling models to autonomously refine their own capabilities.